%----------HEADING-----------------
\begin{tabular*}{\textwidth}{l@{\extracolsep{\fill}}r}
  \textbf{\Large Yufeng (Alex) Xia} & Email: \href{mailto:yufeng.xia@ed.ac.uk}{yufeng.xia@ed.ac.uk} \\
  \href{https://xyf2002.github.io/}{xyf2002.github.io} \quad
  \href{https://www.linkedin.com/in/yufeng-xia-12648a25b}{linkedin.com/in/yufeng-xia-12648a25b} & Mobile: \href{tel:+447918155093}{(+44) 7918155093} \\
  Edinburgh, UK & \\
\end{tabular*}

\vspace{-2pt}

%-----------RESEARCH INTERESTS-----------------
% \section{Research Interests}
%   \begin{itemize}[leftmargin=*]
%     \item\small{I work on \textbf{how high-level computation is turned into efficient parallel execution, and on what that execution costs on real hardware}. This has meant two things. First, \textbf{compiler-driven transformation}: taking a declarative or high-level specification and generating vectorised, quantised or otherwise specialised code automatically, rather than asking a programmer to hand-tune it. Second, \textbf{parallel execution under contention}: deciding how a computation is partitioned across accelerators, and what happens to it when that capacity is taken away mid-execution. I am drawn to problems where the two meet, and would like to bring this experience to the automatic vectorisation and parallelisation of declarative computations and to the systems that execute them at scale.\vspace{-2pt}}
%   \end{itemize}
